\chapter{Fundamentals}
\label{chap:basics}
This chapter introduces the conceptual and technical foundations required to understand the methodology and experiments presented in this thesis. It covers the Reddit platform and its engagement mechanisms, behavioral aspects of online popularity, the machine learning techniques used for textual representation and popularity prediction, and the evaluation metrics used to assess model performance. The purpose of this chapter is to establish common terminology and assumptions, rather than to provide an exhaustive survey of the respective research fields.

\section{Reddit as a Social Media Platform}

% TODO: maybe add one example subreddit with example rules
Reddit is a large-scale social media platform organized into topic-specific communities known as \emph{subreddits}. As shown in \autoref{app:subreddit_rules}, each subreddit defines its own moderation rules, content guidelines, and norms of discourse, resulting in substantial heterogeneity in user behavior across communities. Users can submit posts in multiple formats, including \emph{selftext posts}, which consist of a title and a textual body; \emph{link posts}, which consist of a title and a URL pointing to external content; as well as image and video posts, which center on visual or audiovisual media. However, the focus of this work is limited to text-based content, and therefore only selftext and link posts are considered.

User interaction on Reddit is primarily expressed through upvotes and downvotes. The difference between these votes determines a post’s \emph{score}, which influences its visibility within the subreddit. In addition to score, Reddit exposes an \emph{upvote ratio}, representing the proportion of positive votes, and the number of comments, which serves as a complementary indicator of engagement. These signals capture different aspects of community response and are not interchangeable.

% todo: explain karma  (might use reddiquette website as citation)

The exact mechanics of voting and ranking on Reddit are intentionally opaque. Reddit has historically employed \emph{vote fuzzing}, i.e., the deliberate display of inexact vote totals to hinder spam and vote manipulation \cite{stoddard2015popularity,redditReddiquette,redditJedbergVoteFuzzing}. Reddit’s \textit{Reddiquette}\cite{reddit_reddiquette} still warns that votes are not a reliable basis for polls because of vote fuzzing. An administrator comment linked there further states that the displayed counts are fuzzed for anti-spam reasons and that the discrepancy becomes larger for more active posts \cite{redditReddiquette,redditJedbergVoteFuzzing}. This complicates the interpretation of vote-based engagement signals and motivates caution when treating them as exact quantities. We did not find evidence in the cited sources that the \textit{upvote ratio} is subject to the same fuzzing mechanism.

Importantly, content visibility on Reddit is determined at the subreddit level rather than globally. Ranking algorithms operate within individual communities, meaning that engagement outcomes are inherently context-dependent. A post that is successful in one subreddit may perform poorly in another due to differences in audience size, topical focus, and community norms.\cite{chandrasekharan2018internet, zhang2017community, hamilton2017loyalty} This community-specific nature of engagement is a central consideration throughout this thesis.


\section{Engagement and Popularity on Reddit}

Popularity on Reddit does not correspond to a single observable variable, but instead emerges from a combination of engagement signals such as score, upvote ratio, and comment volume. Prior work has shown that these signals are correlated but capture distinct phenomena, and that high engagement does not necessarily imply consensus or positive reception \cite{hessel2019something}.

Empirical analyses of Reddit and similar platforms indicate that posts typically receive the majority of their engagement shortly after submission, after which attention decays as newer content replaces older posts in ranked listings \cite{terentiev2014predicting,wigsnes2019predicting}. Early visibility therefore plays a critical role in shaping long-term outcomes, reinforcing feedback loops in which initial engagement increases exposure and subsequent interaction.

Beyond temporal dynamics, research in behavioral analytics and computational social science has demonstrated that online popularity is shaped by social reinforcement mechanisms. On Reddit, users respond not only to the content of a post, but also to social signals such as existing votes and comments. This can give rise to herding effects, where users are more likely to engage positively with content that already appears popular \cite{stoddard2015popularity,waller2021quantifying}. At the same time, ideological alignment and community-specific patterns of participation shape which topics and forms of political expression gain traction within a given subreddit \cite{waller2021quantifying}.

These temporal and behavioral mechanisms imply that popularity cannot be fully explained by textual quality alone. Instead, it emerges from the interaction between content, audience expectations, platform affordances, and social dynamics. Understanding these factors is essential for interpreting the performance of predictive models and for assessing the limitations of approaches trained on historical engagement data.

\section{Transformer-Based Text Representations}

Transformer architectures have become the dominant modeling paradigm in natural language processing due to their ability to efficiently capture long-range contextual dependencies in text. Introduced by Vaswani et al.~\cite{Vaswani2017}, the Transformer replaces recurrent and convolutional structures with a self-attention mechanism that allows each token in a sequence to attend to all other tokens in parallel. This design enables scalable training on large corpora and facilitates the learning of rich contextual representations.

\subsubsection{Self-Attention Mechanism}

The self-attention mechanism computes a weighted representation of input tokens by comparing each token against all others. Given an input sequence of embeddings $\mathbf{X} \in \mathbb{R}^{n \times d}$, where $n$ is the sequence length and $d$ is the embedding dimension, self-attention first projects the input into query, key, and value representations:

\begin{align}
\mathbf{Q} &= \mathbf{X}\mathbf{W}_Q \\
\mathbf{K} &= \mathbf{X}\mathbf{W}_K \\
\mathbf{V} &= \mathbf{X}\mathbf{W}_V
\end{align}

where $\mathbf{W}_Q, \mathbf{W}_K, \mathbf{W}_V \in \mathbb{R}^{d \times d_k}$ are learned projection matrices and $d_k$ is the key dimension.

The attention weights are computed using scaled dot-product attention:

\begin{equation}
\text{Attention}(\mathbf{Q}, \mathbf{K}, \mathbf{V}) = \text{softmax}\left(\frac{\mathbf{Q}\mathbf{K}^T}{\sqrt{d_k}}\right)\mathbf{V}
\end{equation}

The scaling factor $\sqrt{d_k}$ prevents the dot products from growing too large in magnitude, which would push the softmax function into regions with vanishingly small gradients~\cite{Vaswani2017}.

Multi-head attention extends this mechanism by computing $h$ parallel attention operations with different learned projections, then concatenating and projecting the results:

\begin{align}
\text{head}_i &= \text{Attention}(\mathbf{Q}\mathbf{W}_Q^i, \mathbf{K}\mathbf{W}_K^i, \mathbf{V}\mathbf{W}_V^i) \\
\text{MultiHead}(\mathbf{Q}, \mathbf{K}, \mathbf{V}) &= \text{Concat}(\text{head}_1, \ldots, \text{head}_h)\mathbf{W}_O
\end{align}

where $\mathbf{W}_O \in \mathbb{R}^{hd_k \times d}$ is the output projection matrix.


Unlike traditional bag-of-words or n-gram models, Transformer-based approaches produce context-sensitive representations in which the meaning of a token depends on its surrounding text. Earlier contextual embedding models, such as ELMo \cite{Peters2018}, demonstrated the importance of context-aware representations, a principle that Transformer architectures generalize and extend. These properties are particularly relevant for social media content, where short texts, informal language, and implicit references are common.

For downstream tasks, Transformer models are often used as encoders that map variable-length text into fixed-dimensional vector representations, commonly referred to as \emph{embeddings}. These embeddings can serve as input features for task-specific models, enabling a separation between representation learning and predictive modeling.

\paragraph{Large Language Models and Text Embeddings}

Large Language Models (LLMs) are Transformer-based models trained on large-scale text corpora using self-supervised learning objectives. While many LLMs are optimized for text generation, a growing class of models is specifically designed to produce high-quality embeddings for downstream tasks. Sentence-BERT~\cite{Reimers2019}, for example, demonstrated that embedding-specialized architectures can outperform generic language models on semantic similarity and clustering tasks.

Recent benchmarks such as the Massive Text Embedding Benchmark (MTEB)\cite{muennighoff2022mteb} provide systematic evaluations of embedding models across a wide range of tasks, including clustering, retrieval, and classification \cite{muennighoff2022mteb}. Results from MTEB indicate that embedding-specialized models often outperform generative LLMs when used purely as feature extractors, while also being more computationally efficient.

A key practical consideration in large-scale embedding-based analyses is dimensionality. Higher-dimensional embeddings can capture more nuanced semantic information, but incur increased storage and computational costs. When encoding millions of documents, these trade-offs become critical. Consequently, embedding models that offer strong performance at moderate dimensionality are preferable for large-scale studies.

In this thesis, Transformer-based embedding models are used exclusively as frozen encoders. No fine-tuning is performed on Reddit data, as the objective is to evaluate how well general-purpose semantic representations transfer to the task of popularity prediction. Among the evaluated models, \textit{stella\_en\_400M\_v5} \cite{stella_en_400M_v5} was selected as a primary candidate due to its competitive performance on semantic benchmarks and its configurable embedding dimensionality, which allows balancing representational capacity and resource constraints.

While embeddings capture semantic properties of text, they do not encode platform-specific dynamics such as ranking algorithms, voting behavior, or temporal effects. As a result, embeddings alone are insufficient to explain engagement outcomes, motivating their combination with additional features and predictive models in subsequent chapters.

\section{Gradient Boosting for Regression}

Gradient Boosting is an ensemble learning technique that constructs a predictive model by iteratively combining multiple weak learners, typically decision trees, into a single strong learner. Rather than training all models independently, Gradient Boosting builds models sequentially, where each new learner is trained to correct the errors made by the ensemble thus far. This is achieved by fitting each learner to the negative gradient of a loss function with respect to the current model predictions, hence the term \emph{gradient} boosting.

\paragraph{Mathematical Formulation}

Gradient boosting constructs an ensemble model $F(\mathbf{x})$ as an additive combination of $M$ base learners $f_m(\mathbf{x})$:

\begin{equation}
F(\mathbf{x}) = \sum_{m=1}^{M} f_m(\mathbf{x})
\end{equation}

The model is trained iteratively by minimizing a differentiable loss function $L(y, F(\mathbf{x}))$. At each iteration $m$, a new base learner $f_m$ is fit to the negative gradient of the loss with respect to the current ensemble prediction:

\begin{equation}
f_m(\mathbf{x}) = \arg\min_{f} \sum_{i=1}^{n} \left( -\frac{\partial L(y_i, F_{m-1}(\mathbf{x}_i))}{\partial F_{m-1}(\mathbf{x}_i)} - f(\mathbf{x}_i) \right)^2
\end{equation}

where $F_{m-1}$ represents the ensemble after $m-1$ iterations.

For regression tasks with squared error loss, the negative gradient corresponds to the residual:

\begin{equation}
-\frac{\partial L(y, F(\mathbf{x}))}{\partial F(\mathbf{x})} = y - F(\mathbf{x})
\end{equation}

XGBoost extends this formulation with an explicit regularization term. The objective function at iteration $m$ becomes:

\begin{equation}
\mathcal{L}^{(m)} = \sum_{i=1}^{n} L(y_i, F_{m-1}(\mathbf{x}_i) + f_m(\mathbf{x}_i)) + \Omega(f_m)
\end{equation}

where the regularization term is defined as:

\begin{equation}
\Omega(f) = \gamma T + \frac{1}{2}\lambda \sum_{j=1}^{T} w_j^2
\end{equation}

Here, $T$ is the number of leaves in the tree, $w_j$ are the leaf weights, $\gamma$ controls tree complexity, and $\lambda$ is the L2 regularization parameter on leaf weights~\cite{chen2016xgboost}.

One of the central advantages of Gradient Boosting is its ability to model complex, non-linear relationships without requiring explicit feature engineering. By combining many shallow decision trees, the model can capture interactions between features while remaining robust to noise and heterogeneous input types. These properties make Gradient Boosting particularly suitable for regression tasks involving a mixture of numerical features and high-dimensional representations.

XGBoost (Extreme Gradient Boosting)\cite{chen2016xgboost} is a widely used and highly optimized implementation of Gradient Boosting that introduces several practical improvements over earlier formulations. These include explicit regularization terms to control model complexity, efficient handling of sparse inputs, and scalable training procedures that allow the model to be applied to large datasets. As a result, XGBoost has become a standard choice for structured data regression tasks and has demonstrated strong performance across a wide range of applications.

In the context of this thesis, popularity prediction is formulated as a regression problem with a continuous target variable representing relative community reception. Gradient Boosting models are well suited for this setting, as they can naturally accommodate non-linear relationships between textual embeddings, temporal features, and auxiliary metadata. Unlike linear regression models, Gradient Boosting does not assume additive or independent effects between features, which is particularly important given the complex and context-dependent nature of social media engagement.

Another practical advantage of XGBoost is its flexibility with respect to feature scaling and distributional assumptions. 
Decision tree–based models are invariant to monotonic transformations of input features and do not require normalization to achieve stable training. 
This property is beneficial when combining dense text embeddings with heterogeneous metadata features.

Finally, the use of Gradient Boosting as a downstream regressor allows for a clear separation between representation learning and prediction. 
Transformer-based models are employed solely to generate fixed text representations, while XGBoost operates on these representations to estimate popularity. 
This modular design simplifies analysis, enables controlled ablation studies, and facilitates comparison between different embedding models without altering the predictive framework.

\section{Evaluation Metrics for Regression Tasks}
\label{sec:evaluation_metrics}

Regression models are typically assessed using multiple complementary metrics that capture different aspects of prediction quality. The following metrics are commonly employed to evaluate continuous predictions against observed outcomes.

\paragraph{Mean Absolute Error (MAE).}

MAE measures the average absolute difference between predicted and observed
upvote ratios. It provides a direct interpretation of prediction error magnitude
within the bounded [0,1] range, since Reddit upvote ratios are naturally
restricted between 0 and 1.

\begin{equation}
MAE = \frac{1}{n} \sum_{i=1}^{n} |y_i - \hat{y}_i|
\end{equation}

where $y_i$ denotes the observed upvote ratio and $\hat{y}_i$ the predicted value.

\paragraph{Root Mean Squared Error (RMSE).}

RMSE measures the square root of the average squared prediction error and
penalizes larger prediction deviations more strongly than MAE.

\begin{equation}
RMSE = \sqrt{\frac{1}{n}\sum_{i=1}^{n}(y_i - \hat{y}_i)^2}
\end{equation}

\paragraph{Coefficient of Determination ($R^2$).} 

The coefficient of determination indicates the proportion of variance in the observed upvote ratios that can be explained by the model predictions:

\begin{equation}
R^2 = 1 - \frac{\sum_{i=1}^{n}(y_i - \hat{y}_i)^2}{\sum_{i=1}^{n}(y_i - \bar{y})^2}
\end{equation}

where $\bar{y} = \frac{1}{n}\sum_{i=1}^{n}y_i$ is the mean of the observed values.

\paragraph{Pearson Correlation Coefficient.} 

Pearson correlation measures the linear relationship between predicted and observed upvote ratios:

\begin{equation}
r = \frac{\sum_{i=1}^{n}(y_i - \bar{y})(\hat{y}_i - \bar{\hat{y}})}{\sqrt{\sum_{i=1}^{n}(y_i - \bar{y})^2}\sqrt{\sum_{i=1}^{n}(\hat{y}_i - \bar{\hat{y}})^2}}
\end{equation}

where $\bar{\hat{y}} = \frac{1}{n}\sum_{i=1}^{n}\hat{y}_i$ is the mean of the predicted values.

Note that while $R^2 = r^2$ holds for ordinary least squares regression, this equality does not apply to non-linear models such as XGBoost used in this study.

Using multiple metrics provides complementary perspectives on model performance. Error-based metrics such as MAE and RMSE evaluate prediction accuracy, while $R^2$ and Pearson correlation measure how well the model captures variability and relative ranking patterns in post popularity.

\section{Statistical Hypothesis Testing}
\label{sec:statistical_testing}

When comparing the performance of different model configurations, it is essential to determine whether observed differences in evaluation metrics reflect genuine effects or merely random variation. Statistical hypothesis testing provides a framework for making such assessments in a principled manner.

\subsection{Wilcoxon Signed-Rank Test}
\label{subsec:wilcoxon}

The Wilcoxon signed-rank test is a non-parametric statistical test used to assess whether the median of paired differences differs significantly from zero~\cite{wilcoxon1945individual}. 
Unlike parametric tests such as the paired t-test, the Wilcoxon signed-rank test does not assume that the data follow a normal distribution, making it appropriate for small sample sizes or data with unknown distributional properties.
The test is commonly applied in paired experimental designs, where each observation unit is measured under two conditions. 
In the context of model evaluation, this corresponds to comparing two configurations on the same set of validation instances, producing paired metric values.

Given $n$ paired observations $(x_1, y_1), \ldots, (x_n, y_n)$, the test procedure is as follows:

\begin{enumerate}
    \item Compute the differences $d_i = y_i - x_i$ for each pair.
    \item Exclude pairs where $d_i = 0$ and adjust $n$ accordingly.
    \item Rank the absolute differences $|d_i|$ from smallest to largest.
    \item Assign each rank the sign of the corresponding difference.
    \item Compute the test statistic as the sum of the positive ranks or the sum of the negative ranks.
\end{enumerate}

Under the null hypothesis that the median difference is zero, the distribution of the test statistic is known and can be used to compute a $p$-value. A small $p$-value provides evidence against the null hypothesis, suggesting that the observed difference is unlikely to have occurred by chance.
The Wilcoxon signed-rank test is particularly well-suited for experiments with a modest number of paired observations, where normality assumptions may not hold but where the paired design allows for effective control of baseline variability.

\subsection{Multiple Comparison Correction}
\label{subsec:multiple_testing}

When conducting multiple hypothesis tests within a single study, the probability of observing at least one false positive increases with the number of tests performed. 
This phenomenon, known as the \emph{multiple testing problem}, can lead to spurious findings if individual $p$-values are interpreted without adjustment.
Several procedures exist to control error rates in multiple testing scenarios. 
One widely used approach is the Benjamini-Hochberg (BH) procedure~\cite{benjamini1995controlling}, which controls the \emph{false discovery rate} (FDR) rather than the more conservative family-wise error rate (FWER).

The false discovery rate is defined as the expected proportion of false positives among all rejected null hypotheses. 
Controlling FDR at level $\alpha$ means that, on average, no more than $\alpha \times 100\%$ of the discoveries are expected to be false positives. 
This approach is less conservative than FWER-controlling methods such as the Bonferroni correction, making it more appropriate when some false positives are acceptable in exchange for greater statistical power.

The Benjamini-Hochberg procedure operates as follows. Given $m$ hypothesis tests with $p$-values $p_1, \ldots, p_m$:

\begin{enumerate}
    \item Order the $p$-values from smallest to largest: $p_{(1)} \leq p_{(2)} \leq \ldots \leq p_{(m)}$.
    \item For a desired FDR level $\alpha$, find the largest index $k$ such that:
    \begin{equation}
    p_{(k)} \leq \frac{k}{m} \alpha
    \end{equation}
    \item Reject the null hypotheses corresponding to $p_{(1)}, \ldots, p_{(k)}$.
\end{enumerate}

Adjusted $p$-values can also be computed to facilitate interpretation. For each test $i$, the adjusted $p$-value is:

\begin{equation}
\tilde{p}_i = \min\left(1, \min_{j \geq i} \left\{ m \cdot \frac{p_{(j)}}{j} \right\} \right)
\end{equation}

These adjusted $p$-values can be compared directly against the desired FDR threshold $\alpha$.

The Benjamini-Hochberg procedure is widely used in contexts where multiple comparisons are inevitable, such as feature selection experiments, hyperparameter sweeps, or ablation studies involving many alternative configurations. By controlling FDR rather than FWER, the procedure balances the need to avoid false discoveries with the goal of detecting true effects.